\documentclass{article}
\usepackage[utf8]{inputenc}
\usepackage[T1]{fontenc}
\usepackage{geometry}
\usepackage{graphicx}
\usepackage{float}
\usepackage{hyperref}
\usepackage[english,indonesian]{babel}

\title{Optimasi Pemrosesan Template LaTeX Menggunakan Pendekatan Berbasis AI}
\author{Academic Support Agent}
\date{\today}

\begin{document}
\maketitle

\begin{abstract}

\end{abstract}

\chapter{PENDAHULUAN}


\chapter{TINJAUAN PUSTAKA}


% Original content (commented out):
% %-------------------------------------------------------------------
% This template was created by Hjalmar Stenberg with updates and
% proofreading by Isak Jonsson. V.1.1 (added support for appendix)
%
% This template is a heavily modified version of the 
% LTU SRT Thesis Template originally created by Johan E. Carlson
% and later updated by Tosin Adewumi.
%
% This template was created due to wanting a simplified version of
% the original template.
%
% Feel free to use and modify this template however you want.
% (may need to reference the original template author idk)
%
% This templete will probably not be maintained and all bugs it may
% possibly have is an exercise for the user to fix.
%
% glhf
%-------------------------------------------------------------------

%-------------------------------------------------------------------
% Document class and package definitions
%-------------------------------------------------------------------
\documentclass[12pt,a4paper]{article}

\input{packages}
\input{glossary}

\begin{document}

%-------------------------------------------------------------------
% Make Titlepage
%-------------------------------------------------------------------
\pagenumbering{roman}
\thispagestyle{empty}

\begin{figure}[h]
 \centering
 \includegraphics[width=0.5\textwidth]{Figures/100 accurate esd.png}
\end{figure}

\begin{center}
    \vspace{1cm} \mbox{}\hrulefill \mbox{}\\[12mm]
    {\Huge \textbf{Thesis template\\[2cm]}}
    {\LARGE \textbf{Hjalmar Stenberg}}\\
    {\large \textbf{feat. Isak Jonsson}}\\[10mm]
    {\large Department of Computer Science, Electrical and Space Engineering\\ 
    Luleå University of Technology\\ 
    Luleå, Sweden} \\[12mm]
    {\large \mbox{}\hrulefill\mbox{}\\[12mm] 
    \textbf{Supervisors:\\[12mm]}}
    Super Viserson, These Sison
\end{center}

\cleardoublepage

%-------------------------------------------------------------------
% Stuff before actual sections. \phantomsection and \clearpage used
% to get hyperrefs in toc working
%-------------------------------------------------------------------
\thispagestyle{plain}
\vspace*{6cm} \hfill{\emph{Whether we wanted it or not...}}

\section*{Abstract}
\input{Sections/0.1 abstract}

\section*{Acknowledgments}
\input{Sections/0.2 preface}

\tableofcontents
\clearpage

\phantomsection
\section*{Acronyms}\label{ch:Acronyms} 
\addcontentsline{toc}{section}{Acronyms}
\printglossary[type=\acronymtype]
\clearpage


%-------------------------------------------------------------------
% If you need to add units to the symbols it can be done by changing 
% to style=long-name-desc-sym in \printglossary and adding
% \renewcommand{\entryname}{Symbol}
% \renewcommand{\symbolname}{Unit}
% Between \addcontentstoline and \printglossary
% This gives correct column header names.
% A change when defining symbols must also be done. Instructions for
% that can be found in glossary.tex 
%-------------------------------------------------------------------
\phantomsection
\section*{List of Symbols}\label{ch:Symbols} 
\addcontentsline{toc}{section}{List of Symbols}
\printglossary[type=symbols,style=long]
\clearpage

\phantomsection
\listoffigures\label{ch:Figures} 
\addcontentsline{toc}{section}{List of Figures}
\clearpage

\phantomsection
\listoftables\label{ch:Tables} 
\addcontentsline{toc}{section}{List of Tables}
\clearpage

%-------------------------------------------------------------------
% Actual sections. More sections can be added by creating new .tex
% files in Sections folder and inputting them here.
%-------------------------------------------------------------------
\section{Introduction}\label{ch:Introduction}
\pagenumbering{arabic}
\input{Sections/1. Introduction}
\clearpage

\section{Background}\label{ch:Background}
\input{Sections/2. Background}
\clearpage

\section{Method}\label{ch:Method}
\input{Sections/3. Method}
\clearpage

\section{Results}\label{ch:Results}
\input{Sections/4. Results}
\clearpage

\section{Discussion}\label{ch:Discussion}
\input{Sections/5. Discussion}
\clearpage

\section{Conclusion}\label{ch:Conclusion}
\input{Sections/6. Conclusion}
\clearpage

\phantomsection
\section*{References}
\addcontentsline{toc}{section}{References}
\printbibliography[heading=none]
\clearpage

%-------------------------------------------------------------------
% Code for handling appendices. More appendices can be created by 
% creating new files and inputting them down below
%-------------------------------------------------------------------
\begin{appendices}

% Appendix section header Format
% Delete if you don't want it to say Appendix X in section title
\titleformat
{\section} % command
[display] % shape
{\normalfont\sc\Large} % format
{} % label
{0.5ex} % sep
{\hrule\medskip\hfill\appendixname\ \thesection: } % before-code
[\medskip\hrule\vspace{5mm}\thispagestyle{plain}] % after-code

% Input your appendices here
\input{Sections/Appendices/1}
\input{Sections/Appendices/2}
\input{Sections/Appendices/3}
\end{appendices}
\end{document}

\end{document}